% Observed quantities are inserted by scripts/build_report.py.
\documentclass[11pt,a4paper]{article}
\usepackage[margin=1in]{geometry}
\usepackage{fontspec}
\setmainfont{Times New Roman}
\setmonofont{Consolas}[Scale=0.9]
\usepackage{amsmath,mathtools}
\usepackage{unicode-math}
\setmathfont{Cambria Math}
\usepackage[expansion=false]{microtype}
\usepackage{graphicx,booktabs,array,tabularx,longtable,caption,enumitem,placeins}
\usepackage{flafter}
\usepackage[hidelinks]{hyperref}
\setlength{\parindent}{1em}
\setlength{\parskip}{0pt}
\setlength{\emergencystretch}{2em}
\clubpenalty=10000
\widowpenalty=10000
\setlength{\tabcolsep}{4pt}
\renewcommand{\arraystretch}{1.1}
\captionsetup{font=small,labelfont=bf}
\renewcommand{\topfraction}{0.9}
\renewcommand{\bottomfraction}{0.8}
\renewcommand{\textfraction}{0.1}
\renewcommand{\floatpagefraction}{0.75}
\setlist{itemsep=0pt,topsep=3pt}
\newcolumntype{Y}{>{\raggedright\arraybackslash}X}
\newcommand{\ResearchFigure}[3]{\begin{figure}[htbp]\centering
\includegraphics[width=\linewidth]{figures/#1.pdf}\caption{#2}\label{#3}\end{figure}}
\title{Experimental Evaluation of Robust Feature Selection}
\author{}
\date{11 October 2026}
\makeatletter
\renewcommand{\maketitle}{\begin{center}{\Large\bfseries\@title\par}\vspace{4pt}
{\small\@date\par}\end{center}\vspace{2pt}}
\makeatother
\begin{document}
\maketitle
\begin{abstract}
We study whether choosing a few variables by minimizing a robust testing loss preserves the difference between two populations. Four synthetic mechanisms distinguish testing power from signal recovery; independent prediction and two real datasets provide further checks. The recorded robust-risk optimum does not consistently give the best testing power or prediction. A separate small-case search comparison reduces full subset evaluations by 24.2\%, while total median-summed time changes by +4.3\%. Fewer inner evaluations therefore do not by themselves establish a speedup.
\end{abstract}

\section{Experimental question and common protocol}\label{sec:design}
Suppose that we may observe only three of twelve variables. Which three should we
choose so that two populations remain distinguishable? Our method chooses a subset
by minimizing the worst-case testing loss over two empirical Wasserstein balls.
The experiments ask whether this choice preserves distribution differences,
recovers the variables responsible for them, and supports prediction on new data.
We then measure the cost of solving the selection problem.

The statistical experiment contains three informative variables and nine noise
variables. The first population has independent standard-normal coordinates.
Only the three informative coordinates change in the second population, as shown
in Table~\ref{tab:mechanisms}. Their positions are shuffled before selection.
Each method chooses three variables from 20 observations per class. Selection,
subsequent classifier fitting, and final testing use disjoint samples; a method
cannot inspect the final observations when choosing its variables.
For each type of difference, the complete selection and evaluation are repeated
30 times with newly generated samples. Figure points show the average result;
intervals reflect uncertainty across those repetitions. Their calculation is
specified in Section~\ref{sec:parameters}.
\begin{table}[htbp]\centering\small
\caption{Four ways in which the informative variables differ. All other variables have identical distributions in both classes.}\label{tab:mechanisms}
\begin{tabularx}{\linewidth}{@{}p{29mm}Y@{}}\toprule
Difference & Change in the second population\\\midrule
Mean & Each informative coordinate has mean $0.9$ instead of zero.\\
Variance & Each has standard deviation $1.8$ instead of one.\\
Dependence & Their pairwise correlation is $0.65$; each individual marginal distribution is unchanged.\\
Mixture & A shared random sign times $1.5$ is added to all three coordinates; the means remain zero.\\\bottomrule
\end{tabularx}\end{table}
\FloatBarrier

These experiments follow the questions studied by Wang et al.~\cite{wang}:
computation, two-sample testing, and variable recovery. They use our own data
mechanisms and fixed sample size and budget. In particular, the dependence case
is an additional difficulty in which individual-coordinate screening has no
population signal. We do not reproduce the original paper's four data models
or its statistical sample-size, dimension, and budget sweeps.

\subsection{What the methods select}
Six practical methods receive the same three-variable budget. The global and
greedy methods optimize the same robust loss; the other methods select variables
using distribution differences, predictive coefficients, or random choice.
\begin{table}[htbp]\centering\small
\caption{Practical selectors. All methods select three variables; implementation parameters are collected in Section~\ref{sec:parameters}.}\label{tab:selectors}
\begin{tabularx}{\linewidth}{@{}p{32mm}Y@{}}\toprule
Selector & How the variables are chosen\\\midrule
Global risk & Search over subsets, using lower bounds to rule out unpromising subsets.\\
Robust greedy & Add one variable at a time, choosing the addition with the smallest fully refitted robust loss.\\
MMD screening & Rank individual coordinates by a kernel measure of the difference between their two distributions.\\
Sparse additive MMD & Jointly choose sparse coordinate weights using the linear additive-kernel criterion of Wang et al.\\
L1 logistic & Rank the coefficient magnitudes of a logistic model with an absolute-value penalty.\\
Random & Draw a three-variable subset uniformly.\\\bottomrule
\end{tabularx}\end{table}
\FloatBarrier

Two additional references help interpret the results. \emph{Known signal} uses the
three true informative variables, which are unavailable to a practical method.
\emph{All features} uses all twelve variables and therefore exceeds the budget.
These references show what happens when the signal is known and when it is kept
together with all the noise. They are not budget-matched competitors.

Our comparison only partially overlaps the baseline study in Wang et al.:
\begin{table}[htbp]\centering\small
\caption{Scope of the comparison with Wang et al. No unimplemented method is claimed as reproduced.}\label{tab:baseline-scope}
\begin{tabularx}{\linewidth}{@{}p{45mm}Y@{}}\toprule
Part of the original study & What is evaluated here\\\midrule
Linear additive-kernel selection & Implemented from the paper's criterion; a test retaining the learned weights is evaluated separately.\\
Quadratic and Gaussian learned kernels & Not implemented.\\
L0-constrained logistic regression & Not implemented. Our L1 coefficient ranking is an additional baseline, not an L0 reproduction.\\
Projected-Wasserstein selection & Not implemented.\\\bottomrule
\end{tabularx}\end{table}
\FloatBarrier

The statistical results below describe subsets produced by the earlier search
implementation. Section~\ref{sec:search} evaluates the revised search separately
on small optimization cases. Those new timings do not establish new statistical
performance or large-scale speed improvements.

\FloatBarrier
\section{Can the selected variables detect a distribution difference?}\label{sec:statistics}
After selection, every subset is tested by the same two-sample procedure. This
comparison evaluates the selected coordinates rather than changing the test
along with the selector. The test uses maximum mean discrepancy (MMD), which
compares within-population kernel similarities with cross-population similarities.
Here the Gaussian radial-basis-function (RBF) kernel is
$k(x,y)=\exp[-\|x-y\|_2^2/(2b^2)]$.

\emph{Power} is the fraction of new batches on which a genuine difference is
detected; larger is better. \emph{False-alarm rate} is the rejection fraction
when both batches come from the same population; it should not systematically exceed the
nominal upper level $0.05$; lower conservative rates are acceptable. We use an unbiased squared-MMD statistic $T$ and a
permutation test:
\begin{equation}\label{eq:permutation}
p=\frac{1+\#\{T_{\rm perm}\ge T_{\rm observed}-10^{-14}\}}{200},\qquad
\text{reject if }p\le0.05.
\end{equation}
The 199 label permutations and final batches are shared across methods.
The kernel bandwidth is fixed using selection data.
Figure~\ref{fig:power} asks how much detectable difference remains after selection. Larger power is better. The two reference supports help separate a poor selection from a difficult test or the dilution caused by noise variables.

For mean differences, global power is 0.928, versus 0.992 for individual-coordinate MMD screening; the paired difference is -6.3 percentage points [-14.3, -1.3] ($n=30$). All 30/30 mean-shift global searches meet the optimization tolerance, yet greedy power is 0.965. Solving the empirical robust objective accurately does not maximize this independent test's power.

\ResearchFigure{power_and_calibration}{Power of the common RBF test: larger is better. The first six methods choose three variables. The last two rows are references: known signal uses the true three-variable support; all features uses twelve variables. Points are seed means and bars are 95\% bootstrap intervals.}{fig:power}
\FloatBarrier
The known-signal powers are 0.957 for variance and 1.000 for mixture differences, whereas using every variable gives 0.138 and 0.207. Keeping all the signal is not enough when irrelevant coordinates weaken the common test. In the dependence case, even the known-signal power is 0.218: the failure of practical selectors is accompanied by limited sensitivity of the 40-observation-per-class test.

A large rejection fraction is useful only if false alarms remain controlled. Table~\ref{tab:null} repeats the common test on two samples from the same population. The nominal upper level is 0.05: false alarms should not systematically exceed it, while lower conservative rates remain acceptable. These are false alarms of the test, not a classifier's class-zero error.
\begin{table}[htbp]\centering\small\setlength{\tabcolsep}{3.5pt}
\caption{False-alarm rate of the common RBF test (declared level 0.05). Entries average testing batches within each training seed, then the 30 seeds. Known signal and all features are diagnostic references.}\label{tab:null}
\begin{tabularx}{\linewidth}{@{}Yrrrr@{}}
\toprule Selector & Mean & Variance & Dependence & Mixture \\
\midrule
Global risk & 0.052 & 0.040 & 0.048 & 0.055 \\
Robust greedy & 0.053 & 0.047 & 0.060 & 0.055 \\
MMD screening & 0.050 & 0.048 & 0.052 & 0.052 \\
Additive MMD & 0.050 & 0.050 & 0.047 & 0.058 \\
L1 logistic & 0.043 & 0.045 & 0.050 & 0.050 \\
Random & 0.048 & 0.050 & 0.047 & 0.052 \\
Known signal & 0.053 & 0.050 & 0.055 & 0.048 \\
All features & 0.055 & 0.057 & 0.050 & 0.058 \\
\bottomrule\end{tabularx}\end{table}
\FloatBarrier
\subsection{What changes when the learned weights are retained?}
The common RBF procedure deliberately discards selector-specific weights. Wang's linear additive method also learns signed coordinate weights. We separately evaluate a test that retains those learned weights. Common power therefore evaluates its chosen support, while native power retains the signed additive weights and evaluates that implemented selector--test combination. Neither column alone is a full replay of the original paper.
\begin{table}[htbp]\centering\small\setlength{\tabcolsep}{3.5pt}
\caption{Two tests on the same Wang-selected support. Common RBF power uses only its coordinates; native power retains the signed additive weights. Entries are seed means and 95\% intervals; native false alarms are also shown.}\label{tab:native}
\begin{tabularx}{\linewidth}{@{}p{24mm}rrrr@{}}
\toprule Difference & $n$ & Common RBF power & Native power & Native false alarms \\
\midrule
Mean & 30 & \shortstack{0.982\\{[0.962, 0.997]}} & \shortstack{0.977\\{[0.950, 0.998]}} & \shortstack{0.045\\{[0.025, 0.065]}} \\
Variance & 30 & \shortstack{0.328\\{[0.222, 0.435]}} & \shortstack{0.518\\{[0.408, 0.625]}} & \shortstack{0.052\\{[0.033, 0.072]}} \\
Dependence & 30 & \shortstack{0.043\\{[0.025, 0.063]}} & \shortstack{0.045\\{[0.030, 0.060]}} & \shortstack{0.052\\{[0.035, 0.070]}} \\
Mixture & 30 & \shortstack{0.763\\{[0.643, 0.872]}} & \shortstack{0.945\\{[0.915, 0.970]}} & \shortstack{0.040\\{[0.025, 0.057]}} \\
\bottomrule\end{tabularx}\end{table}
\FloatBarrier
The native test is substantially stronger in the variance and mixture cases. Treating common RBF power as the performance of the complete additive method would understate that method. The common comparison instead isolates the usefulness of the selected coordinates.


\FloatBarrier
\section{How many irrelevant variables are selected, and how many signals are missed?}\label{sec:recovery}
A test can detect a difference without recovering every informative variable.
We therefore also compare the chosen set $S$ with the true signal set $T_*$.
The false discovery proportion (FDP) is the fraction of selected variables that
are irrelevant. The non-discovery proportion (NDP) is the fraction of true
signals that are missed. Both should be small:
\begin{equation}\label{eq:recovery}
\operatorname{FDP}=\frac{|S\setminus T_*|}{\max\{1,|S|\}},\qquad
\operatorname{NDP}=\frac{|T_*\setminus S|}{|T_*|}.
\end{equation}
For the six practical selectors, $|S|=|T_*|=3$, so FDP and NDP are equal on
each run. A single plotted value reports both, rather than duplicating the same
data. Uniform random selection has expected FDP and NDP $0.75$; perfect recovery
has value zero. The all-feature reference instead has FDP $0.75$ and NDP zero,
because it retains all three signals and all nine noise variables.
Figure~\ref{fig:recovery} now asks whether those useful coordinates are the true informative ones. Smaller FDP/NDP is better. Values near the random expectation 0.75 indicate little recovery beyond uniformly choosing three coordinates.
For mean differences, global FDP/NDP is 0.467, versus 0.244 for MMD screening. For mixture differences, global FDP/NDP is 0.367, versus 0.289 for MMD screening. Thus a high rejection probability can coexist with selecting irrelevant variables: the testing and recovery questions are different.
\ResearchFigure{recovery_and_prediction}{Variable-recovery errors: smaller is better. FDP and NDP coincide for the six three-variable selectors. The lines show perfect recovery (0) and the expected random error (0.75); the observed random selector is a separate point. Points and intervals use training seeds.}{fig:recovery}
\FloatBarrier


\FloatBarrier
\section{Do the selected variables support prediction on new observations?}\label{sec:prediction}
We next hold the prediction model fixed and vary only its selected inputs.
A Gaussian-kernel support-vector machine (SVM) is fitted on a separate sample
for every selected subset. Its balanced error is the mean of the two classwise
error rates; smaller is better, and $0.5$ is chance performance.
Figure~\ref{fig:prediction} compares the same independently fitted SVM on the six practical supports. Lower balanced error means better prediction. This does not reuse the robust selector's fitted test.
For mean differences, errors are 0.315 for global selection and 0.279 for MMD screening. For mixture differences, errors are 0.314 for global selection and 0.289 for MMD screening. The robust selection objective does not consistently produce the best inputs for this different prediction model.
\ResearchFigure{prediction_error}{Independent SVM balanced error: smaller is better and 0.5 is chance performance. All six selectors use the same fitting samples and model settings. Points are seed means and bars are 95\% intervals.}{fig:prediction}
\FloatBarrier
\subsection{How the fitted robust bound compares with fresh-data error}
A separate robust classifier is also refitted on each selected support. Let $p_0(x)$ be its probability of assigning class zero. A randomized decision draws its class from that probability. Its balanced error and the SVM's hard-decision balanced error are
\begin{equation}\label{eq:errors}\begin{aligned}
\operatorname{BER}_{\rm hard}&=\tfrac12(\operatorname{FPR}+\operatorname{FNR}),\\
\operatorname{BER}_{\rm random}&=\tfrac12\{\mathbb E_0[1-p_0(X)]+\mathbb E_1[p_0(Y)]\}.
\end{aligned}\end{equation}
FPR and FNR are the two classwise wrong-decision rates. Unlike a hard classifier, a confident but non-deterministic probability still contributes randomized error. These two error definitions are not interchangeable. Table~\ref{tab:head} compares the robust classifier's fitted bound with its randomized error on new data and under three shifts toward the other class. The shifts have magnitudes 0, 0.25 and 0.50; their maximum is a measured stress error, not a supremum over every distribution in a Wasserstein ball.
\begin{table}[htbp]\centering\small\setlength{\tabcolsep}{3.5pt}
\caption{Refitted robust classifier on globally selected variables. All errors are randomized balanced errors. The fitted bound is the classwise sum-risk upper bound divided by two. Stress is the within-seed maximum over three specified shifts; final columns count paired trials exceeding the fitted bound.}\label{tab:head}
\begin{tabularx}{\linewidth}{@{}p{23mm}rrrrr@{}}
\toprule Difference & $n$ & Fitted bound & Fresh error & Shift maximum & Above bound \\
\midrule
Mean & 30 & 0.251 & 0.448 & 0.462 & 30/30 \\
Variance & 30 & 0.275 & 0.482 & 0.484 & 30/30 \\
Dependence & 30 & 0.287 & 0.497 & 0.501 & 30/30 \\
Mixture & 30 & 0.248 & 0.423 & 0.426 & 30/30 \\
\bottomrule\end{tabularx}\end{table}
\FloatBarrier
Fresh errors exceed the fitted empirical bound in 120/120 paired trials. This is not a certificate failure: the bound concerns the fitted empirical distribution ball, and the unknown population has not been shown to belong to it. A low empirical worst-case loss cannot be read as a demonstrated population-error guarantee.
\subsection{Changing the uncertainty radius}
Table~\ref{tab:radius} changes the radius used both for selection and for refitting the robust classifier. It checks the whole procedure, rather than holding the selected variables fixed. The three radii are preset alternatives, not choices made from the final test outcomes.
\begin{table}[htbp]\centering\small\setlength{\tabcolsep}{3.5pt}
\caption{Joint radius sensitivity. Entries are seed means and 95\% intervals; errors are randomized balanced errors. Stress is the maximum over preset shifts within each seed.}\label{tab:radius}
\begin{tabularx}{\linewidth}{@{}p{22mm}rrrrr@{}}
\toprule Alternative & $\theta$ & $n$ & Power & Fresh error & Shift maximum \\
\midrule
Mean & 0.10 & 10 & \shortstack{0.790\\{[0.595, 0.915]}} & \shortstack{0.494\\{[0.491, 0.497]}} & \shortstack{0.496\\{[0.493, 0.499]}} \\
Mean & 0.25 & 10 & \shortstack{0.855\\{[0.650, 0.980]}} & \shortstack{0.466\\{[0.450, 0.480]}} & \shortstack{0.474\\{[0.461, 0.485]}} \\
Mean & 0.50 & 10 & \shortstack{0.955\\{[0.915, 0.990]}} & \shortstack{0.426\\{[0.411, 0.440]}} & \shortstack{0.442\\{[0.429, 0.453]}} \\
Mixture & 0.10 & 10 & \shortstack{0.495\\{[0.295, 0.710]}} & \shortstack{0.486\\{[0.478, 0.493]}} & \shortstack{0.487\\{[0.480, 0.493]}} \\
Mixture & 0.25 & 10 & \shortstack{0.630\\{[0.425, 0.845]}} & \shortstack{0.446\\{[0.417, 0.469]}} & \shortstack{0.448\\{[0.419, 0.471]}} \\
Mixture & 0.50 & 10 & \shortstack{0.970\\{[0.930, 0.995]}} & \shortstack{0.371\\{[0.352, 0.392]}} & \shortstack{0.378\\{[0.357, 0.401]}} \\
\bottomrule\end{tabularx}\end{table}
\FloatBarrier
Larger radii improve the recorded power and randomized error in these two mechanisms. This is a result of the preset comparison, not a general monotonicity guarantee or an independently validated radius-selection rule.


\FloatBarrier
\section{What happens on real datasets?}\label{sec:real}
WDBC distinguishes benign from malignant breast-cancer observations using
30 available variables; Ionosphere distinguishes good from bad radar returns
using 34 variables~\cite{wdbc,ionosphere}. Each selector keeps three variables.
The subsequent classifiers use disjoint fitting and final observations, as in
the synthetic study. There is no known true signal set, so we report prediction
rather than claim variable-recovery accuracy.
\begin{table}[htbp]\centering\small\setlength{\tabcolsep}{3.5pt}
\caption{Independent common classifiers on real datasets. Entries are mean $\pm$ overlapping-split SD. Robust error is randomized BER; SVM error is hard BER. These are incumbent supports when search is limited.}\label{tab:real}
\begin{tabularx}{\linewidth}{@{}p{18mm}Yrr@{}}
\toprule Data & Selector & \shortstack{Robust\\randomized BER} & SVM BER \\
\midrule
WDBC & Global risk & 0.251\(\pm\)0.046 & 0.081\(\pm\)0.017 \\
WDBC & Robust greedy & 0.220\(\pm\)0.032 & 0.080\(\pm\)0.021 \\
WDBC & MMD screening & 0.179\(\pm\)0.013 & 0.084\(\pm\)0.017 \\
WDBC & Additive MMD & 0.180\(\pm\)0.013 & 0.082\(\pm\)0.014 \\
WDBC & L1 logistic & 0.230\(\pm\)0.034 & 0.078\(\pm\)0.020 \\
WDBC & Random & 0.310\(\pm\)0.066 & 0.156\(\pm\)0.078 \\
Ionosphere & Global risk & 0.300\(\pm\)0.045 & 0.251\(\pm\)0.094 \\
Ionosphere & Robust greedy & 0.285\(\pm\)0.035 & 0.202\(\pm\)0.081 \\
Ionosphere & MMD screening & 0.284\(\pm\)0.044 & 0.195\(\pm\)0.053 \\
Ionosphere & Additive MMD & 0.273\(\pm\)0.042 & 0.181\(\pm\)0.058 \\
Ionosphere & L1 logistic & 0.284\(\pm\)0.034 & 0.198\(\pm\)0.049 \\
Ionosphere & Random & 0.330\(\pm\)0.028 & 0.303\(\pm\)0.069 \\
\bottomrule\end{tabularx}\end{table}
\FloatBarrier
Global search is time-limited in 40/40 splits; its supports must not be called certified optima.


\FloatBarrier
\section{Does reducing full subset evaluations also reduce time?}\label{sec:search}
Evaluating a subset requires solving its continuous robust-testing problem.
We measure both how many such evaluations are needed and the total time spent
choosing and evaluating subsets. The first measures expensive inner work;
the second also includes the integer search used to choose the next subset.
The revision keeps the existing rule for predicting which subset to evaluate next. It additionally retains the stronger dual-fraction lower bound already returned by each full solve, using that information to discard more unpromising subsets. No additional continuous solve is introduced, and every newly chosen binary subset is fully evaluated. We compare it with the original GP-guided tangent search and with refreshing a tangent master at every step.

The separate confirmation study has 16 configurations on 4 generated datasets, with dimensions 7 or 10 and six observations per class. Each method is timed 3 times per configuration. Reported totals sum the median within each configuration; these settings share data and are not 16 independent population trials. Table~\ref{tab:search-confirmation} separates full evaluations from time spent in the integer master.
\begin{table}[htbp]\centering\small\setlength{\tabcolsep}{3.5pt}
\caption{Current confirmation comparison. Counts and seconds sum within-configuration medians; successful-run counts include all three timings. Every full evaluation makes two continuous solves.}\label{tab:search-confirmation}
\begin{tabularx}{\linewidth}{@{}Yrrrr@{}}
\toprule Search & \shortstack{Evaluated\\subsets} & \shortstack{Master\\seconds} & \shortstack{Total\\seconds} & Optimal runs \\
\midrule
Tangent master at every step & 601 & 13.507 & 17.312 & 48/48 \\
Original GP search & 598 & 8.271 & 12.715 & 48/48 \\
Stored fractions & 453 & 9.610 & 13.262 & 48/48 \\
\bottomrule\end{tabularx}\end{table}
\FloatBarrier
The revised search reduces evaluations from 598 to 453, a 24.2\% reduction. Fewer/equal/more evaluations occur in 15/0/1 cases. Total time changes from 12.715 to 13.262 s (+4.3\%), with faster runs in 8/16 cases. The integer master costs more, so saved inner evaluations do not yield an overall speedup here. The unchanged GP cannot be credited as the cause of the evaluation reduction.
Figure~\ref{fig:search-confirmation} shows the paired cases rather than only the totals: fewer evaluations are preferable, while a time ratio above one means slower.
\ResearchFigure{search_confirmation}{Full subset evaluations (upper panel) and revised/original time (lower panel). The dotted line at one means equal time. Cases 1--8 use half hinge and 9--16 logistic loss; within each block, cases are ordered by dimension, seed, and radius pair.}{fig:search-confirmation}
\FloatBarrier
The confirmation solves 144/144 runs to their numerical target; the largest feasible-risk excess over the exhaustive lower reference is \(1.6\times10^{-7}\). Development included one reported master infeasibility; after the common numerical repair, the fixed rule passes 26/26 integration cases. All development comparisons and the failed run are retained in \path{docs/experiment_development_record.md}. The new efficiency evidence is limited to these small cases.


\subsection{Earlier comparison as the dimension increases}\label{sec:compute}
The earlier implementation also compared three price constraints used to link
the binary subset with the continuous problem: individual coordinate caps,
a cap on their weighted sum, and a cap scaled by the risk being minimized.
The last is tested with and without its fractional starting stage.
All four versions solve the same binary robust objective.
Across 159 solved gate/enumeration pairs, the largest recorded feasible-objective difference is 0.000, with 0 incompatible intervals at the combined solver tolerances. At $D\ge20$, global risk reaches the limit in 15/15 mean instances. At $D\ge20$, global risk reaches the limit in 15/15 mixture instances.
\begin{table}[htbp]\centering\small\setlength{\tabcolsep}{3.5pt}
\caption{$D=12$ reference comparison. Time is charged seconds with 95\% seed-bootstrap intervals; supports are distinct integer subsets, excluding fractional root calls.}\label{tab:reference-computation}
\begin{tabularx}{\linewidth}{@{}p{22mm}Yrrr@{}}
\toprule Alternative & Search & Solved & Time (s) & Supports \\
\midrule
Mean & Box & 10/10 & \shortstack{4.07\\{[3.57, 4.88]}} & 220.0 \\
Mean & Triangle & 10/10 & \shortstack{4.63\\{[4.16, 5.08]}} & 219.9 \\
Mean & Risk, no root & 10/10 & \shortstack{13.73\\{[13.05, 14.32]}} & 203.2 \\
Mean & Global risk & 10/10 & \shortstack{16.54\\{[15.74, 17.28]}} & 197.8 \\
Mixture & Box & 10/10 & \shortstack{3.52\\{[3.40, 3.63]}} & 220.0 \\
Mixture & Triangle & 10/10 & \shortstack{4.75\\{[4.38, 5.08]}} & 219.2 \\
Mixture & Risk, no root & 10/10 & \shortstack{13.22\\{[12.30, 14.07]}} & 200.0 \\
Mixture & Global risk & 10/10 & \shortstack{16.39\\{[15.24, 17.52]}} & 197.0 \\
\bottomrule\end{tabularx}\end{table}
\FloatBarrier
\ResearchFigure{optimization_comparison}{Earlier implementation: charged time and solved fraction as dimension grows. Lower time and higher solved fraction are preferable; unsuccessful runs remain included.}{fig:computation}
\FloatBarrier


\FloatBarrier
\section{Implementation parameters and complete records}\label{sec:parameters}
The following settings identify the experiments above. They are fixed before
the final observations are evaluated. The main synthetic comparison uses 30
training seeds; each radius comparison uses 10. Synthetic intervals are 95\%
percentile bootstrap intervals over the seeds of the corresponding experiment,
using 2,000 resamples after averaging test batches within each seed. Real-data entries use mean $\pm$ split standard
deviation because the splits overlap. All time limits and failures remain in
the corresponding denominators and charged computation.
\begin{table}[htbp]\centering\small\setlength{\tabcolsep}{3.5pt}
\caption{Reference statistical data. Selection and classifier fitting counts are per class; final counts show their unit. Repeats are training seeds or overlapping real-data splits.}\label{tab:data}
\begin{tabularx}{\linewidth}{@{}p{25mm}rYrr@{}}
\toprule Data & $D/d$ & \shortstack{Selection / fitting\\per class} & \shortstack{Final\\observations} & Repeats \\
\midrule
Synthetic & 12/3 & 20/20 & 1000/class & 30 \\
WDBC & 30/3 & 20/20 & 489 total & 20 \\
Ionosphere & 34/3 & 20/20 & 271 total & 20 \\
\bottomrule\end{tabularx}\end{table}
\FloatBarrier
\begin{table}[htbp]\centering\small\setlength{\tabcolsep}{3.5pt}
\caption{Reference optimization settings. Each row is evaluated separately for mean and mixture alternatives.}\label{tab:computation-settings}
\begin{tabularx}{\linewidth}{@{}rrrrY@{}}
\toprule $D$ & $d$ & Atoms/class & Seeds & Alternatives \\
\midrule
8 & 2 & 8 & 10 & Mean, mixture \\
12 & 3 & 8 & 10 & Mean, mixture \\
20 & 3 & 10 & 5 & Mean, mixture \\
40 & 3 & 10 & 5 & Mean, mixture \\
80 & 3 & 10 & 5 & Mean, mixture \\
\bottomrule\end{tabularx}\end{table}
\FloatBarrier
The primary class radii are both 0.25. Each synthetic seed has 20 alternative and 20 null batches, with 40 observations per class in a batch. Radius sensitivity uses 10 seeds per mechanism and radius.

Centering and sample-standard-deviation scaling are fitted on pooled selection observations and then frozen. The common RBF bandwidth is their median squared distance. Downstream SVMs use $C=1$ and a median bandwidth fitted only on their own fitting sample. The robust classifier uses the same half-hinge loss and radius as selection.

The sparse additive criterion subtracts $0.1$ times its variance estimate from the weighted kernel mean, with unit-norm sparse weights. Supports are enumerated at $D\le12$; larger cases use forward selection and at most ten improving swaps. We implement the paper's negative variance term, not the positive term found in the inspected script revision, and never refit testing bandwidths using test labels. No cross-validation of this coefficient is performed. L1 logistic ranking uses liblinear, $C=1$, and at most 2,000 iterations. Greedy selection keeps its last completed stage at the time limit and fills remaining slots by MMD screening.

Statistical selection uses 30 s per method, radius 0.25, and numerical gap tolerance $10^{-6}+10^{-5}|\mathrm{UB}|$. The earlier global search starts from MMD screening and uses five fractional starting rounds. Randomized robust predictions use their fitted class probabilities; SVM hard decisions use threshold zero, with ties assigned to class zero.

\begin{table}[htbp]\centering\small\setlength{\tabcolsep}{3.5pt}
\caption{Fixed settings of the separate current-search confirmation.}\label{tab:current-settings}
\begin{tabularx}{\linewidth}{@{}p{36mm}Y@{}}
\toprule Item & Setting \\
\midrule
Cases and data & 16 configurations on 4 datasets; $(D,d)=(7,2),(10,3)$; 6 observations/class; seeds 19,23. \\
Loss and radius & Half hinge or logistic; radii $(0.12,0.20)$ or $(0.30,0.45)$. \\
Generated coordinates & Gaussian class zero; two class-one means shifted by 1.4,0.9; nuisance scales 0.3--3; no fitted standardization. \\
Common starting subset & $|\bar x_{0j}-\bar x_{1j}|/(s_{0j}+s_{1j}+10^{-12})$; sample SD with ddof=0, index ties; no true signal identities. \\
Time and precision & 60 s, one thread; $\varepsilon=2\times10^{-5}$, relative tolerance zero; inner/master targets $\varepsilon/4$. \\
Subset prediction & GP: mean $1/2$, covariance $\frac14 e^{-\|z-z'\|_1/\max(1,d)}$, ridge $\max\{10^{-8},(\varepsilon/(4B))^2\}$. \\
Proposal and refresh & At most 8 proposals; mean minus SD score; master refreshed every 3 full solves or if none remains; no fractional root calls. \\
\bottomrule\end{tabularx}\end{table}
\FloatBarrier
The Gaussian process (GP) predicts normalized risk and its model uncertainty; $B=1$ for half hinge and $B=2\log2$ for logistic normalization. Confirmation counts use the solver's visited-subset counter. Limited historical runs can include an inaccurate final evaluation, while every current confirmation evaluation reaches $\varepsilon/4$. Physical refinements and master time remain in the raw traces. Three timings per configuration rotate execution order.

Reference computations use AMD Ryzen 9 7945HX with Radeon Graphics, one solver thread and SciPy-backed HiGHS 1.12.0; logistic checks use CVXPY 1.9.3 and Clarabel 0.11.1.

The earlier study's 1640 records include 3 failures and 189 time limits. Full outcomes, original attempts and reruns remain in \path{results/research_20261003}. All secondary metrics, supports and optimization traces are available in its \path{per_trial.csv} and JSON records. Development comparisons, integration verification and the earlier fraction comparison are collected in \path{docs/experiment_development_record.md}; current confirmation is stored separately under \path{reports/verification/20261010/next_selection/current}. The report manifests identify the exact inputs used.


\begin{thebibliography}{9}
\bibitem{wang} J. Wang, S. S. Dey, and Y. Xie.
\emph{Variable Selection for Kernel Two-Sample Tests}. 2024.
\href{https://arxiv.org/abs/2302.07415}{arXiv:2302.07415v4}.
Code inspected at commit \texttt{031f740e9ca6} of
\href{https://github.com/WalterBabyRudin/MMDVar_Selection}{MMDVar\_Selection}.
\bibitem{wdbc} UCI Machine Learning Repository.
\emph{Breast Cancer Wisconsin (Diagnostic)}. \href{https://doi.org/10.24432/C5DW2B}{DOI:10.24432/C5DW2B}.
\bibitem{ionosphere} UCI Machine Learning Repository.
\emph{Ionosphere}. \href{https://doi.org/10.24432/C5W01B}{DOI:10.24432/C5W01B}.
\end{thebibliography}
\end{document}
